\documentclass[10pt,a4paper]{amsart}
\usepackage[margin=19mm]{geometry}
\usepackage{amsmath,amssymb,mathtools}
\usepackage[scr=rsfs,cal=euler]{mathalfa}
\usepackage{xfrac}
\usepackage[unicode,hidelinks,bookmarks=false]{hyperref}
\newcommand{\ie}{i.e.\ }
\newcommand{\cf}{cf.\ }
\newcommand{\resp}{respectively\ }
\newcommand{\Subsection}{\subsection}
\providecommand{\nobreakdash}{\nobreak\hbox{-}}
\setlength{\emergencystretch}{2em}
\multlinegap0pt
\usepackage{bbm}
\usepackage{amssymb}
\usepackage{enumerate}
\usepackage{tikz}
\usepackage[shortlabels]{enumitem}
\usepackage{natbib}
\usepackage{caption}
\usepackage[normalem]{ulem}
\newtheorem{thm}{Theorem}
\newcommand{\E}{\mathbb{E}}
\title{Bayesian Quantile Estimation and Regression\\ with Martingale Posteriors}
\author{Edwin Fong, Andrew Yiu}
\date{Source: arXiv:2406.03358v1 (CC BY 4.0)}
\begin{document}
\maketitle

\noindent\emph{Source and licence.} Bayesian Quantile Estimation and Regression\\ with Martingale Posteriors. Version arXiv:2406.03358v1 (CC BY 4.0), distributed by arXiv under the Creative Commons Attribution 4.0 International licence, http://creativecommons.org/licenses/by/4.0/, as recorded in the arXiv licence metadata for this exact version and shown on the abstract page https://arxiv.org/abs/2406.03358v1. Full licence notice: \texttt{licenses/arxiv-2406.03358-CC-BY-4.0.txt}.\par\smallskip

\noindent\textit{Editorial scope.} Counted unit: the article's thm th:tight (source0.tex lines 2445--2447). The article's complete original proof of it is delivered with it (source0.tex lines 2448--2504, one \texttt{proof} environment of 57 lines). No mathematics is written, repaired or re-derived: the statement and the proof are the article's own lines, in the article's own notation, and no retained line is substituted or re-flowed.  Retained uncounted as the source-local context the proof needs: lines 1244--1252.  Nothing was omitted from any retained span, so the package carries no omission record.  No retained line contains a citation, so the wrapper carries no bibliography and no bibliography entry.  No retained line contains a figure environment, an image-inclusion command or drawing code, so no image is delivered and the package declares no asset dependency.  Editorial formatting only, in addition: an AMS \texttt{amsart} preamble that restates the article's own declarations and macros used by the retained lines, namely the environments \texttt{thm} on separate counters, together with the standard packages those lines need.  The retained environments print in this document as Theorem 1 in order of appearance, whereas the article prints the same items with the numbering of its own full document.  The complete native source of the article as distributed in the arXiv e-print for this exact version is bundled unchanged as \texttt{sources/160329/source0.tex}.
\begin{thm}[Bracketing CLT with infinite sums] \label{th::bracketing_clt}
    Suppose that $(\mathcal{F},d)$ is totally bounded and each $Z_{Ni}$ has a finite second moment. Suppose also that
    \begin{align*}
        \sum_{i=N}^{\infty} \E^{*}\Vert Z_{Ni}\Vert_{\mathcal{F}}\mathbbm{1}\{\Vert Z_{Ni}\Vert_{\mathcal{F}} > \eta \} &\rightarrow 0 \quad \text{for every $\eta > 0$,} \\
        \sup_{d(f,g) < \delta_{N}} \sum_{i=N}^{\infty} \E[(Z_{Ni}(f)-Z_{Ni}(g)^{2}] &\rightarrow 0\quad \text{for every $\delta_{N} \downarrow 0$,} \\ 
        \int_{0}^{\delta_{N}} \sqrt{\log \mathfrak{N}_{[]}(\varepsilon,\mathcal{F}, L^{2,N}})\,d\varepsilon &\rightarrow 0\quad \text{for every $\delta_{N} \downarrow 0$.} 
    \end{align*}
    Then the sequence $\sum_{i=N}^{\infty}(Z_{Ni}-\E Z_{Ni})$ is asymptotically tight in $\ell^{\infty}(\mathcal{F})$ and converges in distribution provided it converges marginally. If the partitions can be chosen independent of $N$, then the middle of the displayed conditions is unnecessary.
\end{thm}

\begin{thm}\label{th:tight}
The sequence of functions $\sqrt{N}S_N$ is  asymptotically tight in $\ell^{\infty}(\mathcal{F})$ with probability 1.
\end{thm}

\begin{proof}
    We verify the assumptions in {Theorem \ref{th::bracketing_clt}}. Our semimetric space is $(\mathcal{F} = (0,1)\times \{1,\ldots, p\}, d)$, where $d((u_{1},j_{1}),(u_{2},j_{2})) = |u_{1}-u_{2}|^{1/2} + \mathbbm{1}(j_{1} \neq j_{2})$. We have used the discrete metric on $\{1,\ldots, p\}$ and then specified the sum of the two semimetrics to define the semimetric product space. Clearly, this semimetric space is totally bounded.

First define $Z_{Ni}(u, j) = \sqrt{N}\alpha_{i}(H_{\rho_{i}}(u,V_{i})-u)X_{ij}$ for all $N$ and $i \geq N$. Note that this definition differs in nature to the non-regression case (where the $-u$ term can be omitted) because the randomness in the covariates must be accounted for. A trivial envelope function for $Z_{Ni}$ is $F_{Ni} = 2\sqrt{N}\alpha_{i}C_{X}$. We need to verify the Lindeberg condition
\begin{equation*}
    \sum_{i=N}^{\infty} F_{Ni}\mathbbm{1}\{F_{Ni} > \eta\} \rightarrow 0
\end{equation*}
for every $\eta > 0$. Since $F_{Ni} < 2cC_{X}/\sqrt{N}$, we will have $F_{Ni} < \eta$ for all sufficiently large $N$. Thus, the Lindeberg condition holds.

Next we need 
\begin{equation*}
    \sup_{d((u_{1},j_{1}),(u_{2},j_{2})) < \delta_{N}} \sum_{i=N}^{\infty} \E\left\{Z_{Ni}(u_{1})-Z_{Ni}(u_{2})\right\}^{2} \rightarrow 0
\end{equation*}
for every $\delta_{N} \downarrow 0$. For all sufficiently large $N$, we must have $\delta_{N} < 1$, in which case
\begin{equation*}
    d((u_{1},j_{1}),(u_{2},j_{2})) < \delta_{N} \implies j_{1} = j_{2} \quad \text{and} \quad |u_{1} - u_{2}|^{1/2} < \delta_{N}.
\end{equation*}
Let $u_{1} > u_{2}$ with $|u_{1} - u_{2}|^{1/2}  < \delta_{N}$. We have
\begin{align*}
     \E\left\{Z_{Ni}(u_{1}, j)-Z_{Ni}(u_{2},j)\right\}^{2}&\leq  2\E\left\{\sqrt{N}\alpha_{i}X_{ij}H_{\rho_{i}}(u_{1},V_{i})-\sqrt{N}\alpha_{i}X_{ij}H_{\rho_{i}}(u_{2},V_{i})\right\}^{2} \\
     &+ 2\E\left\{\sqrt{N}\alpha_{i}X_{ij}u_{1}-\sqrt{N}\alpha_{i}X_{ij}u_{2}\right\}^{2}\\
     &\leq  2N\alpha_{i}^{2}C_{X}^{2}\E\left\{H_{\rho_{i}}(u_{1},V_{i})-H_{\rho_{i}}(u_{2},V_{i})\right\}^{2} \\
     &+ 2N\alpha_{i}^{2}C_{X}^{2}\delta_{N}^{2}
\end{align*}

Note that $H_{\rho_{i}}(u,V_{i})$ is non-decreasing in $u$ with $H_{\rho_{i}}(0,V_{i}) = 0$ and $H_{\rho_{i}}(1,V_{i}) = 1$. So
\begin{align*}
    \E\left\{H_{\rho_{i}}(u_{1},V_{i})-H_{\rho_{i}}(u_{2},V_{i})\right\}^{2} &\leq \E\left\{H_{\rho_{i}}(u_{1},V_{i})-H_{\rho_{i}}(u_{2},V_{i})\right\} \\
    &= (u_{1}-u_{2})\\
    &< \delta_{N}^{2}.
\end{align*}
Thus, for all sufficiently large $N$ such that $\delta_{N} < 1$, we have
\begin{equation*}
    \sup_{d((u_{1},j_{1}),(u_{2},j_{2})) <\delta_{N}} \sum_{i=N}^{\infty}\E\left\{Z_{Ni}(u_{1}, j_{1})-Z_{Ni}(u_{2}, j_{2})\right\}^{2} < 4N\delta_{N}^{2} C_{X}^{2}\sum_{i=N}^{\infty} \alpha_{i}^{2}.
\end{equation*}
Since $\limsup_{N \rightarrow \infty} N \sum_{i=N}^{\infty} \alpha_{i}^{2} \leq c^{2}$, the right-hand side of the above display tends to zero for any $\delta_{N} \downarrow 0$.

Finally, we need to verify the bracketing entropy integral condition. Fix $j \in \{1,\ldots, p\}$ for the time being. For all sufficiently large $n$, we will have $N \sum_{i=N}^{\infty} \alpha_{i}^{2} \leq 2c^{2}$. Given $\varepsilon > 0$, choose a partition $0 = u_{0} < u_{1} < \ldots < u_{M} = 1$ such that $u_{k} - u_{k-1} < \varepsilon^{2}/(8C_{X}^{2}c^{2})$ for every $k$. The number of points in the partition can be chosen to be smaller than a constant times $1/\varepsilon^{2}$. Then
\begin{align*}
    \E \sup_{u_{k-1} \leq s,t < u_{k}} |Z_{Ni}(s,j) - Z_{Ni}(t,j)|^{2} &< 2N\alpha^{2}_{i}C_{X}^{2}\E\left\{H_{\rho_{i}}(u_{k},V_{i})-H_{\rho_{i}}(u_{k-1},V_{i})\right\}^{2}\\
    &+ 2N\alpha^{2}_{i}C_{X}^{2}\E \sup_{u_{k-1} \leq s,t < u_{k}} |s-t|^{2} \\
    &< 2N\alpha^{2}_{i}C_{X}^{2} \E\left\{H_{\rho_{i}}(u_{k},V_{i})-H_{\rho_{i}}(u_{k-1},V_{i})\right\}\\
    &+ 2N\alpha^{2}_{i}C_{X}^{2}(u_{k} - u_{k-1}) \\
    &< \frac{N \alpha_{i}^{2}\varepsilon^{2}}{2c^{2}}.
\end{align*}
We deduce that for all sufficiently large $N$,
\begin{equation*}
    \sum_{i=N}^{\infty}\E \sup_{u_{j-1} \leq s,t < u_{j}} |Z_{Ni}(s) - Z_{Ni}(t)|^{2} < \varepsilon^{2}.
\end{equation*}
In other words, $\mathfrak{N}_{[]}(\varepsilon, \mathcal{F}, L_{2,n}) \lesssim \frac{1}{\varepsilon^{2}}$, which verifies the entropy condition. 

Thus, {Theorem \ref{th::bracketing_clt}} implies that
\begin{equation*}
    -\sum_{i=N}^{\infty} Z_{Ni}(u,j) = \sqrt{N}\sum_{i=N}^{\infty}\alpha_{i}(u - H_{\rho_{i}}(u,V_{i}))X_{ij}
\end{equation*}
is asymptotically tight in $\ell^{\infty}(\mathcal{F})$.
\end{proof}

\end{document}
